% Options for packages loaded elsewhere
\PassOptionsToPackage{unicode}{hyperref}
\PassOptionsToPackage{hyphens}{url}
\documentclass[
  a4paper,
]{article}
\usepackage{xcolor}
\usepackage[margin=25mm]{geometry}
\usepackage{amsmath,amssymb}
\setcounter{secnumdepth}{-\maxdimen} % remove section numbering
\usepackage{iftex}
\ifPDFTeX
  \usepackage[T1]{fontenc}
  \usepackage[utf8]{inputenc}
  \usepackage{textcomp} % provide euro and other symbols
\else % if luatex or xetex
  \usepackage{unicode-math} % this also loads fontspec
  \defaultfontfeatures{Scale=MatchLowercase}
  \defaultfontfeatures[\rmfamily]{Ligatures=TeX,Scale=1}
\fi
\usepackage{lmodern}
\ifPDFTeX\else
  % xetex/luatex font selection
\fi
% Use upquote if available, for straight quotes in verbatim environments
\IfFileExists{upquote.sty}{\usepackage{upquote}}{}
\IfFileExists{microtype.sty}{% use microtype if available
  \usepackage[]{microtype}
  \UseMicrotypeSet[protrusion]{basicmath} % disable protrusion for tt fonts
}{}
\makeatletter
\@ifundefined{KOMAClassName}{% if non-KOMA class
  \IfFileExists{parskip.sty}{%
    \usepackage{parskip}
  }{% else
    \setlength{\parindent}{0pt}
    \setlength{\parskip}{6pt plus 2pt minus 1pt}}
}{% if KOMA class
  \KOMAoptions{parskip=half}}
\makeatother
\setlength{\emergencystretch}{3em} % prevent overfull lines
\providecommand{\tightlist}{%
  \setlength{\itemsep}{0pt}\setlength{\parskip}{0pt}}
\usepackage{bookmark}
\IfFileExists{xurl.sty}{\usepackage{xurl}}{} % add URL line breaks if available
\urlstyle{same}
\hypersetup{
  pdftitle={Emergent Spacetime and Wave-Scale Readout from Anonymous Complex Zero Closure: Minimal Readout{,} an Infinite Constructive Solution{,} Logarithmic Spacetime Mapping{,} and Post-Mapping \textbackslash lambda\textbackslash nu Invariance},
  pdfauthor={Noriaki Kihara (Independent Researcher), ORCID 0009-0004-6753-4020},
  hidelinks,
  pdfcreator={LaTeX via pandoc}}

\title{Emergent Spacetime and Wave-Scale Readout from Anonymous Complex
Zero Closure: Minimal Readout, an Infinite Constructive Solution,
Logarithmic Spacetime Mapping, and Post-Mapping \(\lambda\nu\)
Invariance}
\author{Noriaki Kihara (Independent Researcher), ORCID
0009-0004-6753-4020}
\date{September 3, 2026 --- v2.0 --- Version DOI
10.5281/zenodo.22282218}

\begin{document}
\maketitle

\subsection{Abstract}\label{abstract}

This paper does not place space, time, wavelength, frequency, or
velocity as foundational variables in advance. Starting only from an
anonymous set of complex numbers identified under permutations and the
complex square-zero closure

\[
\sum_n z_n^2=0,
\]

we investigate, in order of minimal cardinality, what can be read out
from the inside.

First, a single complex number formally carries the quadratic form

\[
z^2=(a^2-b^2)+2iab,
\]

but we distinguish this from the fact that, once the common complex
scale is quotiented out as a gauge, no nontrivial relative quantity can
be obtained from one element alone. For two nonzero complex numbers,
imposing

\[
z_1^2+z_2^2=0
\]

yields

\[
z_2=\pm i z_1,
\]

so that the amplitude ratio \[1\] and the relative phase \[\pm\pi/2\]
are fixed without any external measuring rod. For three elements, the
squared quantities form a closed triangle in the complex plane,
producing a cosine-rule-type constraint between relative amplitude
ratios and relative phases. In general a continuous degree of freedom
remains here, and we do not claim that a finite phase lattice arises
automatically from square-zero closure alone.

Next, as the central construction of this paper, we give a countably
infinite special solution built from an integer sequence \[M_k\] and
phase cycles within each shell,

\[
z_{k,n}
=
\frac{A}{M_k}
\exp\left[i\left(\phi_k+\frac{\pi n}{M_k}\right)\right],
\qquad
n=0,\ldots,M_k-1.
\]

Every shell closes exactly under the square-zero condition, and
moreover, if

\[
\sum_k \frac{1}{M_k}<\infty,
\]

the infinite sum of squares converges absolutely. In particular
\[M_k=k(k+1)\] is an explicit example that satisfies absolute
convergence and \[M_{k+1}/M_k\to1\] simultaneously.

In this special solution, while anonymity is preserved, a relative
spatial scale

\[
L_k\propto M_k
\]

can be read from the shell amplitudes, and a relative period scale

\[
P_k=M_k
\]

can be read from the cycle order of the squared phase set. Therefore,
before the mapping, \[L_k/P_k\] is constant, and setting
\[\nu_k^{(0)}:=1/P_k\] gives \[L_k\nu_k^{(0)}=\mathrm{const}\].

Furthermore, since the continuous homomorphism that carries
multiplicative scale ratios to additive relative coordinates is the
logarithm up to a constant factor, we define, between two shells
\[j,k\],

\[
\Delta X_{jk}=C_X\ln\frac{L_j}{L_k},
\qquad
\Delta T_{jk}=C_T\ln\frac{P_j}{P_k}.
\]

After this spacetime readout, setting

\[
\lambda_{jk}:=|\Delta X_{jk}|,
\qquad
\tau_{jk}:=|\Delta T_{jk}|,
\qquad
\nu_{jk}:=\frac{1}{\tau_{jk}},
\]

the special solution's \[L_k\propto P_k\] again yields

\[
\lambda_{jk}\nu_{jk}
=
\frac{C_X}{C_T}
=
\mathrm{const}.
\]

Therefore, in this construction, the \[\lambda\nu\]-type reciprocal
relation is preserved not only before the spacetime mapping but also
after the actual mapping to additive spatial and temporal scales.

Finally, combining the two readouts, converted to the same units, into a
single complex readout

\[
W=X+iT,
\]

the original squaring operation gives

\[
dW^2=dX^2-dT^2+2i\,dX\,dT,
\]

whose real part is a \[1+1\]-dimensional Lorentzian quadratic form. This
paper does not claim that this construction is a unique derivation of
real spacetime. What it shows is that the anonymous complex zero-closure
system admits an explicit infinite special solution from which relative
angles, relative scales, period scales, additive spacetime scales, a
\[\lambda\nu\]-type invariant, and a Lorentzian signature can be read
out in one consistent order.

\textbf{Keywords}: anonymous complex numbers, square-zero closure,
readout, emergent spacetime, relative scale, logarithmic mapping,
wavelength, frequency, Lorentzian metric, roots of unity

\begin{center}\rule{0.5\linewidth}{0.5pt}\end{center}

\subsection{1. Problem Setting --- Not Values, but ``What Can Be
Read''}\label{problem-setting-not-values-but-what-can-be-read}

In ordinary wave theory, one first places the spatial coordinate \[x\],
the time \[t\], the wavelength \[\lambda\], and the frequency \[\nu\],
and then describes a wave as

\[
\exp\{i(kx-\omega t)\}.
\]

This paper reverses the order.

What is given first is only a set of complex numbers and the closure of
the sum of squares:

\[
\mathcal Z_N
=
\{z_1,\ldots,z_N\}/S_N,
\]

\[
Q(\mathcal Z_N)
=
\sum_{n=1}^{N}z_n^2
=0.
\tag{1}
\]

The quotient by \[S_N\] is taken so that no complex number is given a
physical name or order. Computational indices are mere memory locations,
and readouts must not depend on permutations.

Equation (1) is also preserved under the common transformation by any
\[\gamma\in\mathbb C^\times\],

\[
z_n\mapsto \gamma z_n,
\]

under which

\[
Q\mapsto \gamma^2Q=0.
\]

Therefore equation (1) alone does not fix the common amplitude and the
common phase. In this paper, quantities that depend on this common
degree of freedom are not admitted as absolute readouts; differences,
ratios, closure orders, and permutation-invariant symmetries are the
candidates for internal readouts.

The question of this paper is simple.

\begin{quote}
Starting from an anonymous set of complex numbers with no background
space and no background time, and importing no external measuring rod,
at which stage can one construct relative quantities corresponding to
angle, scale, period, space, time, wavelength, and frequency?
\end{quote}

\begin{center}\rule{0.5\linewidth}{0.5pt}\end{center}

\subsection{2. One Complex Number --- Distinguishing the Formal Square
Readout from the Gauge-Invariant
Readout}\label{one-complex-number-distinguishing-the-formal-square-readout-from-the-gauge-invariant-readout}

Let a single complex number be

\[
z=a+ib.
\]

Its square is

\[
z^2=(a^2-b^2)+2iab.
\tag{2}
\]

Therefore, with respect to the fixed complex algebraic basis \[(1,i)\],

\begin{itemize}
\tightlist
\item
  the difference quadratic form \[a^2-b^2\]
\item
  the cross term \[2ab\]
\end{itemize}

exist formally inside a single complex number.

However, if the common complex scale \[z\mapsto\gamma z\], and in
particular the common phase rotation, is quotiented out as a gauge, then
the individual values of \[a\] and \[b\], and the numerical value of
\[ab\], are in general not invariant. This paper therefore distinguishes

\begin{enumerate}
\def\labelenumi{\arabic{enumi}.}
\tightlist
\item
  \textbf{that the square form exists algebraically}
\item
  \textbf{that it can be read as a relative quantity even after the
  gauge is quotiented out}.
\end{enumerate}

With a single element there is nothing to compare with; under the full
\[\mathbb C^\times\] gauge, all nonzero points belong to a single orbit.
The strict sense of relative readout adopted in this paper therefore
begins with two elements.

This distinction becomes important later. The Lorentzian quadratic form
reappears at the final stage from the same square structure as equation
(2), but at that point \[X,T\] are quantities already constructed as
relative readouts.

\begin{center}\rule{0.5\linewidth}{0.5pt}\end{center}

\subsection{3. Two Elements --- The First Relative Angle without an
External Angular
Gauge}\label{two-elements-the-first-relative-angle-without-an-external-angular-gauge}

Suppose two nonzero elements \[z_1,z_2\] satisfy

\[
z_1^2+z_2^2=0.
\tag{3}
\]

Then

\[
z_2^2=-z_1^2,
\]

and

\[
\left(\frac{z_2}{z_1}\right)^2=-1.
\]

Therefore

\[
\boxed{
\frac{z_2}{z_1}=\pm i
}
\tag{4}
\]

holds.

Equation (4) immediately gives

\[
\boxed{
\frac{|z_2|}{|z_1|}=1
}
\tag{5}
\]

and

\[
\boxed{
\arg z_2-\arg z_1
\equiv\frac{\pi}{2}
\pmod{\pi}
}
\tag{6}
\]

Here the absolute phase is not determined. Rotating the two elements
simultaneously preserves equation (3). However, \textbf{the right angle
as a relative phase} does not disappear under common rotation.

Therefore the nontrivial zero closure of two elements is the minimal
configuration that fixes the relative amplitude \[1\] and the relative
angle \[\pi/2\] without first placing an external protractor.

Note that the \[\pi/2\] obtained here must not be confused with a claim
that all phases of a general many-body system are quantized to four
values. What equation (3) fixes is the relative angle intrinsic to the
two-element closure.

\begin{center}\rule{0.5\linewidth}{0.5pt}\end{center}

\subsection{4. Three Elements --- Coupling of Relative Amplitudes and
Relative
Phases}\label{three-elements-coupling-of-relative-amplitudes-and-relative-phases}

Next consider three nonzero elements with

\[
x^2+y^2+z^2=0.
\tag{7}
\]

Writing the squared quantities as

\[
X=x^2,
\qquad
Y=y^2,
\qquad
Z=z^2,
\]

we have

\[
X+Y+Z=0.
\tag{8}
\]

That is, \[X,Y,Z\] form a closed triangle in the complex plane. Since
\[Z=-(X+Y)\],

\[
|Z|^2
=
|X|^2+|Y|^2
+2|X||Y|\cos(\arg X-\arg Y).
\]

Using \[|X|=|x|^2\] and \[\arg X=2\arg x\], we obtain

\[
\boxed{
\cos 2(\arg x-\arg y)
=
\frac{|z|^4-|x|^4-|y|^4}
{2|x|^2|y|^2}
}
\tag{9}
\]

The right-hand side of equation (9) consists only of relative
amplitudes; the left-hand side consists only of relative phases. Thus,
for the first time with three elements, the amplitude ratio, which
collapsed to \[1\] for two elements, acquires a nontrivial shape degree
of freedom, and this shape and the relative phase are bound by the same
closure condition.

In this paper we call this the \textbf{conjugate-type amplitude--phase
readout}. The word ``conjugate-type'' is used in the sense that fixing
one constrains the other; it does not mean that canonically conjugate
variables have been proved.

Also, equation (9) in general admits continuous shape parameters.
Therefore, from

\[
\sum_{j=1}^{3}z_j^2=0
\]

alone, one cannot say that relative phases are automatically restricted
to a finite lattice such as \[45^\circ\] or \[90^\circ\]. Finite phase
cycles are \textbf{explicitly realized as the structure of the solution}
in the infinite special solution below.

\begin{center}\rule{0.5\linewidth}{0.5pt}\end{center}

\subsection{5. A Constructive Special Solution with Countably Infinite
Zero
Closure}\label{a-constructive-special-solution-with-countably-infinite-zero-closure}

\subsubsection{5.1 General shell
construction}\label{general-shell-construction}

Take an integer sequence

\[
M_k\in\mathbb N,
\qquad
M_k\ge2,
\qquad
k=1,2,\ldots
\]

With \[A\in\mathbb C^\times\] a common scale and \[\phi_k\in\mathbb R\]
a common phase for each shell, place \[M_k\] complex numbers for each
\[k\]:

\[
\boxed{
z_{k,n}
=
\frac{A}{M_k}
\exp\left[i\left(\phi_k+\frac{\pi n}{M_k}\right)\right]
}
\qquad
(n=0,\ldots,M_k-1)
\tag{10}
\]

Squaring gives

\[
z_{k,n}^2
=
\frac{A^2e^{2i\phi_k}}{M_k^2}
\exp\left(\frac{2\pi i n}{M_k}\right).
\tag{11}
\]

Therefore, from the sum of \[M_k\]-th roots of unity,

\[
\sum_{n=0}^{M_k-1}z_{k,n}^2
=
\frac{A^2e^{2i\phi_k}}{M_k^2}
\sum_{n=0}^{M_k-1}
\exp\left(\frac{2\pi i n}{M_k}\right)
=0.
\tag{12}
\]

Every shell closes exactly under the square-zero condition.

\subsubsection{5.2 Absolute convergence of the infinite
sum}\label{absolute-convergence-of-the-infinite-sum}

The absolute value of each squared term is

\[
|z_{k,n}^2|
=
\frac{|A|^2}{M_k^2},
\]

so for the whole system

\[
\sum_{k=1}^{\infty}
\sum_{n=0}^{M_k-1}|z_{k,n}^2|
=
|A|^2
\sum_{k=1}^{\infty}\frac{1}{M_k}.
\tag{13}
\]

Hence, if

\[
\boxed{
\sum_{k=1}^{\infty}\frac{1}{M_k}<\infty
}
\tag{14}
\]

the sum of squares converges absolutely. Therefore the sum does not
change under any rearrangement of the terms as an anonymous set, and

\[
\boxed{
\sum_{k,n}z_{k,n}^2=0
}
\tag{15}
\]

holds exactly.

\subsubsection{5.3 The minimal explicit
example}\label{the-minimal-explicit-example}

In this paper we take

\[
\boxed{
M_k=k(k+1)
}
\tag{16}
\]

as the representative example. Then

\[
\sum_{k=1}^{\infty}\frac{1}{M_k}
=
\sum_{k=1}^{\infty}
\frac{1}{k(k+1)}
=1.
\tag{17}
\]

Furthermore,

\[
\frac{M_{k+1}}{M_k}
=
\frac{(k+1)(k+2)}{k(k+1)}
=
\frac{k+2}{k}
\longrightarrow1.
\tag{18}
\]

Therefore this sequence simultaneously satisfies

\begin{enumerate}
\def\labelenumi{\arabic{enumi}.}
\tightlist
\item
  absolute convergence of the infinite sum of squares,
\item
  shell cycle order \[M_k\to\infty\],
\item
  adjacent scale ratio \[M_{k+1}/M_k\to1\].
\end{enumerate}

\begin{center}\rule{0.5\linewidth}{0.5pt}\end{center}

\subsection{6. Re-identification of Shells from the Anonymous
Set}\label{re-identification-of-shells-from-the-anonymous-set}

In each shell of equation (10), the amplitudes of all elements coincide:

\[
|z_{k,n}|
=
\frac{|A|}{M_k}.
\tag{19}
\]

If the \[M_k\] are chosen to be mutually distinct, the shells can be
reconstructed from the anonymous set as input, through the
\textbf{multiplicity of equal amplitudes}.

That is, letting \[N(r_k)\] be the number of elements with amplitude
\[r_k\], the special solution has

\[
N(r_k)=M_k,
\qquad
r_k=\frac{|A|}{M_k}.
\tag{20}
\]

Therefore

\[
\boxed{
r_kN(r_k)=|A|
}
\tag{21}
\]

is common to all shells.

Even though \[|A|\] itself is a common scale and cannot be read, the
ratio between shells \[j,k\],

\[
\frac{r_j}{r_k}
=
\frac{M_k}{M_j},
\tag{22}
\]

can be read.

On the other hand, the squared phase set of each shell is

\[
\left\{
\exp\left(2i\phi_k+\frac{2\pi i n}{M_k}\right)
\right\}_{n=0}^{M_k-1},
\]

which forms a regular \[M_k\]-gon. Its minimal cyclic generator is

\[
U_k
=
\exp\left(\frac{2\pi i}{M_k}\right),
\]

\[
U_k^{M_k}=1.
\tag{23}
\]

Therefore \[M_k\] is reconstructed identically both from the amplitude
multiplicity and from the phase-cycle order.

That ``two independent readouts return the same integer \[M_k\]'' is the
central structure of this special solution.

\begin{center}\rule{0.5\linewidth}{0.5pt}\end{center}

\subsection{\texorpdfstring{7. Pre-Mapping Spatial and Period Scales and
the \(\lambda\nu\)
Readout}{7. Pre-Mapping Spatial and Period Scales and the \textbackslash lambda\textbackslash nu Readout}}\label{pre-mapping-spatial-and-period-scales-and-the-lambdanu-readout}

\subsubsection{7.1 Relative spatial scale}\label{relative-spatial-scale}

Absolute amplitudes are not read. However, the ratio of inverse
amplitudes is

\[
\frac{|z_k|^{-1}}{|z_j|^{-1}}
=
\frac{M_k}{M_j}.
\]

We therefore read the space-type relative scale \[L_k\] obtained from
amplitudes, up to a common constant, as

\[
\boxed{
L_k\propto M_k
}
\tag{24}
\]

This is not a declaration of an absolute length such as ``\[M_k\]
meters''. What can be read is only the ratio

\[
\boxed{
\frac{L_j}{L_k}=\frac{M_j}{M_k}
}
\tag{25}
\]

\subsubsection{7.2 Relative period scale}\label{relative-period-scale}

The squared phases of the same shell return to the start after repeating
the minimal rotation \[2\pi/M_k\] exactly \[M_k\] times. We therefore
take the period-type readout to be

\[
\boxed{
P_k:=M_k
}
\tag{26}
\]

No second has been introduced here either. \[P_k\] is ``the number of
internal phase steps required for one full cycle''.

Its inverse is defined as the frequency-type readout,

\[
\boxed{
\nu_k^{(0)}:=\frac{1}{P_k}
=\frac{1}{M_k}
}
\tag{27}
\]

\subsubsection{7.3 The pre-mapping reciprocal
relation}\label{the-pre-mapping-reciprocal-relation}

From equations (24) and (27), up to unit normalization,

\[
\boxed{
L_k\nu_k^{(0)}
=\mathrm{const}
}
\tag{28}
\]

holds.

However, equation (28) alone is not sufficient. \[L_k\] and \[P_k\] are
still multiplicative scale labels, and one must verify whether the same
relation is preserved after mapping to ordinary additive spatial and
temporal coordinates.

This is done in the next section.

\begin{center}\rule{0.5\linewidth}{0.5pt}\end{center}

\subsection{8. Logarithmic Readout from Multiplicative Scales to
Additive
Spacetime}\label{logarithmic-readout-from-multiplicative-scales-to-additive-spacetime}

\subsubsection{8.1 Why the logarithm}\label{why-the-logarithm}

Relative scale ratios are multiplicative. For three shells \[i,j,k\],

\[
\frac{L_k}{L_i}
=
\frac{L_k}{L_j}\frac{L_j}{L_i}.
\tag{29}
\]

A continuous function \[f\] carrying this to additive relative
coordinate differences

\[
X_{ki}=X_{kj}+X_{ji}
\]

must satisfy

\[
f(ab)=f(a)+f(b).
\]

Assuming continuity on the positive reals, its general form is

\[
\boxed{
f(r)=C\ln r
}
\tag{30}
\]

Therefore the logarithm is not an arbitrarily chosen convenient
transformation; it is \textbf{unique up to a constant factor as the
continuous homomorphism that turns multiplicative scale ratios into
additive relative distances}.

\subsubsection{8.2 Spatial readout}\label{spatial-readout}

Define the space-type relative interval between two shells \[j,k\] as

\[
\boxed{
\Delta X_{jk}
:=
C_X\ln\frac{L_j}{L_k}
}
\tag{31}
\]

In the special solution \[L_k\propto M_k\], so

\[
\Delta X_{jk}
=
C_X\ln\frac{M_j}{M_k}.
\tag{32}
\]

No absolute origin is needed, and \[X\mapsto X+X_0\] does not change the
readout. In this sense \[X\] is a relative affine coordinate in one
direction.

\subsubsection{8.3 Time-scale readout}\label{time-scale-readout}

Ratios of periods are also multiplicative:

\[
\frac{P_k}{P_i}
=
\frac{P_k}{P_j}\frac{P_j}{P_i}.
\]

Therefore, by the same logic carrying period scales to additive
time-scale differences,

\[
\boxed{
\Delta T_{jk}
:=
C_T\ln\frac{P_j}{P_k}
}
\tag{33}
\]

In the special solution \[P_k=M_k\], so

\[
\Delta T_{jk}
=
C_T\ln\frac{M_j}{M_k}.
\tag{34}
\]

A caution is needed here. What is directly obtained from the static
anonymous set is the cyclic symmetry \[C_{M_k}\] of each shell and the
period scale \[P_k\]. Equation (33) is a \textbf{relative time scale
comparing the period scales of different clocks}; it is not a derivation
of the arrow of time or of an infinite winding number. The phase states
inside a shell give an \[S^1\]-type clock phase, but a complete time
order \[T\in\mathbb R\] requires additional dynamics or an unwrap rule.

With this limitation made explicit, this paper examines the consistency
of the relative metric between spatial and temporal scales.

\begin{center}\rule{0.5\linewidth}{0.5pt}\end{center}

\subsection{\texorpdfstring{9. \(\lambda\nu=\mathrm{const}\) after the
Spacetime
Mapping}{9. \textbackslash lambda\textbackslash nu=\textbackslash mathrm\{const\} after the Spacetime Mapping}}\label{lambdanumathrmconst-after-the-spacetime-mapping}

This is the main verification of this paper.

Read the post-mapping spatial scale between two shells as

\[
\boxed{
\lambda_{jk}
:=
|\Delta X_{jk}|
}
\tag{35}
\]

Similarly, take the post-mapping period scale as

\[
\boxed{
\tau_{jk}
:=
|\Delta T_{jk}|
}
\tag{36}
\]

and define its inverse as

\[
\boxed{
\nu_{jk}
:=
\frac{1}{\tau_{jk}}
}
\tag{37}
\]

In the special solution \[L_j/L_k=P_j/P_k=M_j/M_k\], so

\[
\lambda_{jk}
=
C_X\left|\ln\frac{M_j}{M_k}\right|,
\]

\[
\tau_{jk}
=
C_T\left|\ln\frac{M_j}{M_k}\right|.
\]

Hence

\[
\boxed{
\lambda_{jk}\nu_{jk}
=
\frac{C_X}{C_T}
=:\,c_*
}
\tag{38}
\]

Taking \[C_X=C_T\] in normalized units gives

\[
\boxed{
\lambda_{jk}\nu_{jk}=1.
}
\tag{39}
\]

When returning to dimensional units, \[c_*=C_X/C_T\] becomes the common
constant converting spatial scales and temporal scales.

What matters is that equation (39) is not concluded merely from the
pre-mapping identity

\[
M_k\times\frac{1}{M_k}=1.
\]

Rather, \textbf{after} actually mapping

\[
M
\longmapsto
X=C_X\ln M,
\]

\[
M=P
\longmapsto
T=C_T\ln P
\]

to additive spacetime scales, we confirm that the \[\lambda\] and
\[\nu\] redefined from intervals on that spacetime satisfy equation
(38).

Therefore, in this special solution, the \[\lambda\nu\]-type relation is
not destroyed by the spacetime readout; it commutes with the readout
mapping.

\begin{center}\rule{0.5\linewidth}{0.5pt}\end{center}

\subsection{10. The Lorentzian Quadratic
Form}\label{the-lorentzian-quadratic-form}

To convert the spatial and temporal scales to the same units, set

\[
\widetilde T
:=
\frac{C_X}{C_T}T.
\]

Below we again write \[T\] for simplicity of notation.

Combine the relative readouts into a single complex quantity

\[
\boxed{
W=X+iT
}
\tag{40}
\]

Applying the same ``complex squaring'' as the starting point of this
paper gives

\[
\boxed{
dW^2
=dX^2-dT^2+2i\,dX\,dT
}
\tag{41}
\]

Therefore the real part is

\[
\boxed{
\operatorname{Re}(dW^2)
=dX^2-dT^2
}
\tag{42}
\]

which carries the \[1+1\]-dimensional Lorentzian signature \[(+,-)\].

The imaginary part is

\[
\boxed{
\operatorname{Im}(dW^2)
=2\,dX\,dT
}
\tag{43}
\]

which is of the same form as the cross term appearing in the square of a
single complex number,

\[
(a+ib)^2=a^2-b^2+2iab.
\]

Furthermore, in the special solution, for any two shells \[j,k\] the
same scale ratio is read on the spatial side and on the temporal side,
so

\[
\frac{\Delta X_{jk}}{C_X}
=
\frac{\Delta T_{jk}}{C_T}
=
\ln\frac{M_j}{M_k}
\]

holds. Here \[C_X,C_T\] are the constants defining the continuous
homomorphisms from multiplicative scales to additive coordinates, and
their signs correspond to the orientations of the respective readout
axes. Since the underlying construction does not fix a positive
direction of the axes, normalizing to the same absolute units
\[|C_X|=|C_T|\] allows us to write

\[
\Delta X_{jk}
=
\sigma\,\Delta T_{jk},
\qquad
\sigma=\pm1.
\]

Therefore, on equation (42), for every shell pair,

\[
\boxed{
ds_{jk}^2
=
\Delta X_{jk}^2-\Delta T_{jk}^2
=0
}
\tag{44}
\]

holds. That is, in this special solution it is not that some null-type
branch exists: \textbf{every shell-pair interval is null}.

This is another expression of the same structure as
\[\lambda_{jk}\nu_{jk}=c_*\] in Section 9. In this special solution,

\[
L_k\propto P_k
\]

simultaneously yields the post-mapping

\[
\lambda\nu=\mathrm{const}
\]

and, in identical units,

\[
ds^2=0.
\]

Therefore the self-duality between ``the spatial scale read from
amplitudes'' and ``the period scale read from phase cycles'', the
\[\lambda\nu\]-type invariant, and the null readout are three
expressions of one identity structure in this special solution.

What is claimed here is not that all physics of real Lorentzian
spacetime has been derived. Equation (42) is the algebraic result that a
Lorentzian quadratic form naturally appears when the two relative scales
constructed in this paper are fed back into the complex square readout.
Also, the fact that all shell pairs become null is a property of this
special solution; it does not mean that general square-zero-closure
solutions are restricted to null readouts. Full Lorentz symmetry, causal
order, the physical realization of boosts, and the extension to \[3+1\]
dimensions remain separate problems.

\begin{center}\rule{0.5\linewidth}{0.5pt}\end{center}

\subsection{11. The Asymptotic Continuum
Limit}\label{the-asymptotic-continuum-limit}

For the representative example \[M_k=k(k+1)\], the adjacent differences
of the spatial and temporal scales are

\[
\Delta X_k
=
C_X\ln\frac{M_{k+1}}{M_k}
=
C_X\ln\frac{k+2}{k}
\longrightarrow0,
\tag{45}
\]

\[
\Delta T_k
=
C_T\ln\frac{k+2}{k}
\longrightarrow0.
\tag{46}
\]

This refinement is not uniform; it becomes exponentially finer depending
on the readout coordinate. Fixing a reference and taking

\[
X_k=C_X\ln M_k
=C_X\ln[k(k+1)],
\]

we have for large \[k\]

\[
X_k\sim2C_X\ln k,
\qquad
k\sim\exp\left(\frac{X_k}{2C_X}\right).
\]

On the other hand, equation (45) gives

\[
\Delta X_k
=C_X\ln\left(1+\frac{2}{k}\right)
\sim\frac{2C_X}{k},
\]

so the local resolution at coordinate \[X\] is asymptotically

\[
\boxed{
\Delta X(X)
\sim
2C_X
\exp\left(-\frac{X}{2C_X}\right)
}
\tag{46a}
\]

Therefore the continuum limit of this construction does not assume a
uniform real lattice from the start; it is a \textbf{position-dependent
asymptotic continuization} in which the discrete steps are exponentially
refined as the logarithmic readout coordinate advances.

Also, the minimal step of the squared phases of each shell is

\[
\Delta\theta_k
=
\frac{2\pi}{M_k}
=
\frac{2\pi}{k(k+1)}
\longrightarrow0.
\tag{47}
\]

Therefore, for large \[k\],

\begin{enumerate}
\def\labelenumi{\arabic{enumi}.}
\tightlist
\item
  the adjacent spacing of the spatial scale axis,
\item
  the adjacent spacing of the temporal scale axis,
\item
  the step of the clock phase
\end{enumerate}

all become arbitrarily fine simultaneously.

In this paper we call this the \textbf{asymptotic local continuum
limit}. The discrete set does not become the real line itself on a
finite region from the start, but since the relative resolution goes to
zero as \[k\to\infty\], it has a local limit approximating a continuous
metric.

In general, for integer sequences with the same properties in this
construction it suffices to require at least

\[
M_k\to\infty,
\]

\[
\sum_k\frac{1}{M_k}<\infty,
\]

\[
\frac{M_{k+1}}{M_k}\to1.
\tag{48}
\]

For example, if \[M_k\] is an integerized \[k^p\] type with \[p>1\], one
obtains a wide family of special solutions satisfying these three
conditions.

\begin{center}\rule{0.5\linewidth}{0.5pt}\end{center}

\subsection{12. Generalization of the
Construction}\label{generalization-of-the-construction}

The essence of equation (10) is not the particular \[k(k+1)\]. It is to
place

\[
M_k
\]

phase-cycle points in shell \[k\] and to tie their common amplitude to

\[
|z_k|
\propto
\frac{1}{M_k}.
\tag{49}
\]

By this condition, the same integer \[M_k\] appears both in

\begin{itemize}
\tightlist
\item
  the inverse scale obtained from amplitudes, and
\item
  the cycle order obtained from the phase set.
\end{itemize}

The generalized sufficient condition is therefore

\[
\boxed{
\text{inverse-amplitude scale}
\ \propto\
\text{phase-cycle order}
}
\tag{50}
\]

If equation (50) holds, the multiplicative relative spatial scale and
the relative period scale have the same ratio, so that even after the
logarithmic mapping

\[
\lambda\nu=c_*
\]

is preserved.

This is a sufficient condition for isomorphic constructions beyond the
special solution of this paper. However, this paper does not prove that
equation (50) is necessary, nor that equation (50) necessarily emerges
from general solutions of anonymous zero closure.

\begin{center}\rule{0.5\linewidth}{0.5pt}\end{center}

\subsection{13. Relation to Prior Work}\label{relation-to-prior-work}

Each mathematical component of this paper has known nearby research.
However, the starting point differs from the order of readout adopted
here.

Penrose's twistor theory is the representative research describing
Minkowski spacetime and conformal geometry using complex projective
geometry and null structures {[}1{]}. There, however, Minkowski
spacetime, or its complexification and conformal structure, is already
positioned as the theoretical background. This paper differs in its
problem setting: it inputs no spacetime coordinates and constructs
relative spacetime scales from the internal readout of anonymous complex
zero closure.

Nottale's scale relativity introduces relativistic structure into scale
transformations themselves and discusses Lorentz-type scale
transformations using logarithmic scales and a generalization of the de
Broglie relations {[}2{]}. The closeness to this paper lies in treating
multiplicative scales in logarithmic coordinates. In contrast, this
paper does not give physical lengths or spacetime resolutions from the
start; it first reconstructs the comparative scales themselves from the
amplitude multiplicities and phase-cycle orders of an anonymous complex
set.

Englman and Yahalom derived, from analyticity, an exact reciprocal
relation between the phase and the log modulus for time-dependent wave
functions {[}3{]}. That work is an important precedent showing that the
phase and the logarithmic amplitude of a complex quantity are not
independent. However, in that work the time \[t\] is given from the
start as the independent variable of the wave function \[\psi(t)\]. The
aim of this paper is to read out period scales and spatial scales from
an anonymous set without background time.

In the literature survey at the time of writing, no prior work was found
that executes the chain

\[
\text{anonymous complex set}
\to
\sum z_n^2=0
\to
\text{minimal relative readout}
\to
\text{explicit infinite zero-closure solution}
\to
(L,P)
\to
(X,T)
\to
\lambda\nu=\mathrm{const}
\to
\operatorname{Re}(dW^2)=dX^2-dT^2
\]

from the same starting point. However, this is not a complete proof of
nonexistence in the literature; it carries the qualification \textbf{to
the best of our knowledge}.

\begin{center}\rule{0.5\linewidth}{0.5pt}\end{center}

\subsection{14. What Is Shown as Theorems and What Remains as Physical
Interpretation}\label{what-is-shown-as-theorems-and-what-remains-as-physical-interpretation}

To avoid overstating the claims of this paper, we separate what is shown
directly in mathematics from what is adopted as readout interpretation.

\subsubsection{14.1 Shown directly}\label{shown-directly}

\begin{enumerate}
\def\labelenumi{\arabic{enumi}.}
\tightlist
\item
  From the nontrivial two-element closure \[z_1^2+z_2^2=0\] follows
  \[z_2/z_1=\pm i\].
\item
  From the three-element closure follows the amplitude--phase constraint
  of equation (9).
\item
  The shells of equation (10) close exactly under the square-zero
  condition for each \[k\].
\item
  Under the condition \[\sum1/M_k<\infty\], the infinite sum of squares
  converges absolutely.
\item
  The same \[M_k\] can be read both from the multiplicity of shell
  amplitudes and from the cycle order of squared phases.
\item
  The continuous homomorphisms from the multiplicative group
  \[\mathbb R_{>0}\] to the additive group \[\mathbb R\] are constant
  multiples of \[\ln\].
\item
  If \[L_k\propto P_k\], then \[\lambda\nu=\mathrm{const}\] of equation
  (38) holds also after the logarithmic spacetime mapping.
\item
  The real part of the square of \[W=X+iT\] is \[dX^2-dT^2\].
\end{enumerate}

\subsubsection{14.2 Adopted as readout
interpretation}\label{adopted-as-readout-interpretation}

\begin{enumerate}
\def\labelenumi{\arabic{enumi}.}
\tightlist
\item
  Calling the inverse-amplitude ratio a space-type scale \[L\].
\item
  Calling the phase-cycle order a period-type scale \[P\].
\item
  Interpreting \[\ln L\] and \[\ln P\] as relative spatial and relative
  temporal scales, respectively.
\item
  Reading the post-mapping spatial scale difference as a wavelength
  \[\lambda\] and the inverse of the temporal scale difference as a
  frequency \[\nu\].
\item
  Reading \[\operatorname{Re}(dW^2)\] as a candidate physical spacetime
  line element.
\end{enumerate}

These are mutually consistent and close on the same special solution.
However, their uniqueness, and their identity with real physics, are
subjects of future verification.

\begin{center}\rule{0.5\linewidth}{0.5pt}\end{center}

\subsection{15. Limitations and Next Verification
Problems}\label{limitations-and-next-verification-problems}

The problems left unsolved by this paper are clear.

\subsubsection{15.1 Not a general
solution}\label{not-a-general-solution}

Equation (10) is a constructive special solution with strong symmetry.
It is not shown that general solutions of the anonymous complex zero
closure

\[
\sum_n z_n^2=0
\]

necessarily have shell structure or equation (50).

\subsubsection{15.2 Reducible shell by
shell}\label{reducible-shell-by-shell}

Since each shell closes to zero on its own, the whole system is
reducible as a direct sum of local closures. The stronger next problem
is whether the same readout holds for irreducible infinite closures that
cannot be decomposed into finite subsystems.

\subsubsection{15.3 The arrow of time is not
derived}\label{the-arrow-of-time-is-not-derived}

What is obtained from the static set is phase cycles and period scales.
The physical process unwrapping the clock phase from \[S^1\] to
\[\mathbb R\], the time order, and the causal direction are not derived
in this paper.

\subsubsection{\texorpdfstring{15.4 The extension from \(1+1\) to
\(3+1\) is
unresolved}{15.4 The extension from 1+1 to 3+1 is unresolved}}\label{the-extension-from-11-to-31-is-unresolved}

What is constructed in this paper is one spatial scale direction and one
temporal scale direction. By which constraints three independent spatial
directions are selected from an anonymous complex set is a separate
problem.

\subsubsection{15.5 No absolute values of physical constants are
produced}\label{no-absolute-values-of-physical-constants-are-produced}

The \[c_*\] of equation (38) is a conversion constant between two
readout units. This paper does not identify its numerical value with the
known speed of light \[c\]. Connecting to real wave propagation requires
independent observation or dynamical verification.

These limitations do not damage the constructive existence proof of this
paper. Rather, they sharpen the next falsifiable question:

\begin{quote}
In more general or irreducible anonymous zero closures as well, do the
amplitude-derived relative scale and the phase-derived period scale
factorize through the same internal parameter, and is the \[\lambda\nu\]
invariance after the spacetime mapping still preserved?
\end{quote}

\begin{center}\rule{0.5\linewidth}{0.5pt}\end{center}

\subsection{16. Conclusion}\label{conclusion}

This paper did not start from ``what is spacetime''; it started from
``what can be read without an external measuring rod''.

The two-element complex zero closure fixes, through

\[
z_2=\pm iz_1,
\]

the first gauge-invariant relative amplitude and relative angle. In the
three-element closure, the relative amplitudes and relative phases are
bound within the same constraint through the closed triangle in the
squared plane.

Furthermore, constructing the countably infinite special solution

\[
z_{k,n}
=
\frac{A}{M_k}
\exp\left[i\left(\phi_k+\frac{\pi n}{M_k}\right)\right],
\]

exact zero closure holds in each shell, and under \[\sum1/M_k<\infty\]
the infinite sum also converges absolutely. In this set, while anonymity
is preserved, the double readout

\[
\text{inverse amplitude}
\longleftrightarrow M_k
\longleftrightarrow
\text{phase-cycle order}
\]

holds.

As a result, before the mapping, the space-type scale \[L_k\] and the
period-type scale \[P_k\] satisfy

\[
L_k\propto P_k,
\]

and mapping both from multiplicative to additive scales by the logarithm
gives

\[
\Delta X
=C_X\ln\frac{L_j}{L_k},
\qquad
\Delta T
=C_T\ln\frac{P_j}{P_k}.
\]

The most important result is that for the wavelength- and frequency-type
quantities redefined \textbf{after this spacetime readout},

\[
\boxed{
\lambda\nu=c_*
}
\]

is preserved.

Furthermore, re-applying the complex square readout to

\[
W=X+iT
\]

yields the Lorentzian quadratic form

\[
\boxed{
\operatorname{Re}(dW^2)
=dX^2-dT^2
}
\]

Therefore this paper gives, for the anonymous complex zero-closure
system, one closed construction path:

\[
\boxed{
\text{minimal relative readout}
\to
\text{amplitude--phase coupling}
\to
\text{infinite closure special solution}
\to
\lambda\nu\text{-type readout}
\to
(X,T)\text{ readout}
\to
\text{recovery of }\lambda\nu
\to
\text{Lorentzian signature}
}
\]

This is not a claim to have uniquely derived real spacetime. The central
result of this paper is to have shown that within complex zero closure
with no background spacetime, there exists at least one infinite special
solution from which relations isomorphic to the ordinary wave and
spacetime metric can be read out self-consistently to the end.

\begin{center}\rule{0.5\linewidth}{0.5pt}\end{center}

\subsection{References}\label{references}

{[}1{]} R. Penrose, ``Twistor Algebra,'' \emph{Journal of Mathematical
Physics}, \textbf{8}(2), 345--366 (1967). DOI:
\href{https://doi.org/10.1063/1.1705200}{10.1063/1.1705200}.

{[}2{]} L. Nottale, ``The Theory of Scale Relativity,''
\emph{International Journal of Modern Physics A}, \textbf{7}(20),
4899--4936 (1992). DOI:
\href{https://doi.org/10.1142/S0217751X92002222}{10.1142/S0217751X92002222}.

{[}3{]} R. Englman and A. Yahalom, ``Reciprocity between moduli and
phases in time-dependent wave functions,'' \emph{Physical Review A},
\textbf{60}, 1802--1810 (1999). DOI:
\href{https://doi.org/10.1103/PhysRevA.60.1802}{10.1103/PhysRevA.60.1802}.

\begin{center}\rule{0.5\linewidth}{0.5pt}\end{center}

\subsection{Appendix A: Uniqueness of the Logarithmic
Readout}\label{appendix-a-uniqueness-of-the-logarithmic-readout}

Suppose a map \[f\] carrying positive scale ratios \[r>0\] to additive
coordinates satisfies

\[
f(rs)=f(r)+f(s).
\tag{A1}
\]

Setting \[r=e^u\] and defining

\[
g(u):=f(e^u),
\]

we get

\[
g(u+v)=g(u)+g(v).
\tag{A2}
\]

If continuity is required of \[f\], hence of \[g\], the only continuous
solutions of Cauchy's additive equation are

\[
g(u)=Cu.
\]

Hence

\[
\boxed{
f(r)=C\ln r.
}
\tag{A3}
\]

Therefore the logarithm used for \[X,T\] in this paper is unique up to a
constant factor under the requirement of carrying the multiplication of
scale ratios to the addition of coordinate differences.

\begin{center}\rule{0.5\linewidth}{0.5pt}\end{center}

\subsection{\texorpdfstring{Appendix B: The Three Conditions for the
Representative Sequence
\(M_k=k(k+1)\)}{Appendix B: The Three Conditions for the Representative Sequence M\_k=k(k+1)}}\label{appendix-b-the-three-conditions-for-the-representative-sequence-m_kkk1}

For the representative sequence

\[
M_k=k(k+1),
\]

\subsubsection{B.1 Absolute convergence}\label{b.1-absolute-convergence}

\[
\sum_{k=1}^{\infty}\frac{1}{M_k}
=
\sum_{k=1}^{\infty}
\left(\frac1k-\frac1{k+1}\right)
=1.
\]

\subsubsection{B.2 Unbounded resolution}\label{b.2-unbounded-resolution}

\[
M_k\to\infty.
\]

Therefore the phase step within a shell

\[
\frac{2\pi}{M_k}\to0.
\]

\subsubsection{B.3 Local continuization of relative
scales}\label{b.3-local-continuization-of-relative-scales}

\[
\frac{M_{k+1}}{M_k}
=1+\frac{2}{k}
\to1,
\]

therefore

\[
\ln\frac{M_{k+1}}{M_k}
\to0.
\]

From the above, \[M_k=k(k+1)\] is one of the simplest explicit examples
that simultaneously satisfies the absolute convergence of the infinite
closure and the asymptotic refinement of the spatial, temporal, and
phase scales.

\end{document}
